\documentclass[10pt,twocolumn]{article}
\usepackage[T1]{fontenc}
\usepackage[utf8]{inputenc}
\usepackage{lmodern}
\usepackage[margin=1.8cm,top=2.4cm,bottom=2cm,columnsep=0.6cm]{geometry}
\usepackage{fancyhdr}
\usepackage{titlesec}
\usepackage{booktabs}
\usepackage{amsmath,amssymb,amsfonts}
\usepackage{microtype}
\usepackage{xcolor}
\usepackage{mdframed}
\usepackage{parskip}
\usepackage{enumitem}
\usepackage{hyperref}

\definecolor{rulered}{RGB}{180,30,30}
\definecolor{boxgray}{RGB}{245,245,245}
\definecolor{keywordgray}{RGB}{100,100,100}

\pagestyle{fancy}
\fancyhf{}
\fancyhead[L]{\textsc{\small Claude Mag}}
\fancyhead[C]{\small\textit{Machine Learning Research Digest}}
\fancyhead[R]{\small 2026-09-09}
\fancyfoot[C]{\small\thepage}
\renewcommand{\headrulewidth}{0.8pt}

\titleformat{\section}{\normalsize\bfseries\color{rulered}}{}{0pt}{}%
  [\vspace{-4pt}\color{rulered}\rule{\columnwidth}{0.4pt}]
\titlespacing{\section}{0pt}{8pt}{4pt}

\hypersetup{colorlinks=true,urlcolor=rulered,linkcolor=black}

\begin{document}
\setlength{\parindent}{0pt}
\setlength{\parskip}{4pt}

{\small\color{keywordgray}\textbf{Keywords:}
  KV cache \textbullet{}
  speculative eviction \textbullet{}
  efficient reasoning \textbullet{}
  LLM inference}\par\vspace{4pt}

{\LARGE\bfseries\raggedright
  Random Attention: Rethinking KV Cache Eviction
  for Efficient Reasoning}\par\vspace{4pt}

{\large\itshape\raggedright
  Scoring Cached Tokens Adds Nothing ---
  Random Eviction with Prompt Protection Matches
  the Best Evictors at 32--43\% Higher Throughput}\par\vspace{6pt}

{\small\textbf{By} Heng Wang, Jielin Qiu, Wenting Zhao, Cheng Qian,
  Liangwei Yang, Jiawei Han, Heng Ji, Silvio Savarese,
  Shelby Heinecke, Huan Wang
  (Salesforce AI Research; UIUC; Cornell; et al.)}\par\vspace{2pt}

{\small\color{keywordgray}arXiv:2609.03430 \textbullet{} EMNLP 2026}
\par\vspace{2pt}\noindent{\color{rulered}\rule{\columnwidth}{0.6pt}}\vspace{6pt}

Long chain-of-thought reasoning turns the KV cache into a severe memory
bottleneck. Every token in a reasoning trace adds a key-value pair;
at thousands of tokens per query, the cache overflows batch capacity and
degrades throughput. Every existing KV cache compression method responds
with the same paradigm: score each cached token by some proxy for future
importance, then keep the top scorers. Random Attention discards the
scores entirely. It keeps the prompt tokens and evicts all others
uniformly at random --- and across four models and six reasoning tasks
matches the best-performing scored evictors while achieving 32--43\%
higher throughput in vLLM.

\begin{mdframed}[backgroundcolor=boxgray,linecolor=rulered,linewidth=1pt,
    innerleftmargin=6pt,innerrightmargin=6pt,
    innertopmargin=6pt,innerbottommargin=6pt]
\textbf{\small AT A GLANCE}\par\vspace{4pt}
\begin{description}[leftmargin=0pt,labelsep=4pt,itemsep=2pt,parsep=0pt,
    font=\small\bfseries]
  \item[Problem:] Long CoT reasoning makes the KV cache a memory
    bottleneck; existing eviction methods rely on expensive importance
    scoring.
  \item[Why it matters:] Production inference systems serving reasoning
    models need maximal throughput; scoring overhead is non-trivial.
  \item[Method:] Keep prompt KV cache; evict reasoning-trace tokens
    uniformly at random per attention head. No score computed.
  \item[Result:] Matches strongest evictors on 6 tasks; 32--43\%
    higher throughput at 50\% cache budget in vLLM.
  \item[Evidence:] Ablations isolating prompt protection from scoring;
    redundancy analysis across text and attention heads.
\end{description}
\end{mdframed}

\section{The Central Observation}

Every KV cache evictor in the literature shares a core design assumption:
the identity of which tokens to retain matters, and a good scoring
function is needed to identify them. H2O uses cumulative attention
weight; StreamingLLM uses recency with attention sinks; LESS uses
gradient-based estimates; SnapKV clusters similar keys.

Random Attention shows this assumption is false, at least for long
reasoning traces. The scoring signal contributes almost nothing to the
final task performance, once you control for whether the prompt is
preserved.

\section{Why the Scoring Signal Is Worthless}

The authors run controlled ablation experiments across all four model
families:

\textbf{Condition A:} Prior evictor, prompt not explicitly protected
--- the prompt may happen to be evicted.

\textbf{Condition B:} Prior evictor with explicit prompt protection.

\textbf{Condition C:} Random eviction, no prompt protection.

\textbf{Condition D (Random Attention):} Random eviction with prompt
protection.

The results show Condition~B $\approx$ Condition~D $\gg$
Conditions~A~and~C. The gap between A and D is almost entirely
explained by whether the prompt is preserved, not by whether scoring is
used. The gap between B and D (evictor scoring vs.\ random after prompt
protection) is negligible.

\section{Two Mechanisms Explaining Redundancy}

Why does random eviction of the reasoning trace lose almost nothing?

\textbf{Text-level redundancy.} Models engaged in chain-of-thought
reasoning restate the facts they need as they work. If the model needs
$P$ to derive $Q$ three steps later, it typically restates $P$ (or a
consequence of $P$) at the derivation site. Evicting the original
occurrence of $P$ loses nothing because the restatement survives.
The authors measure text-level redundancy: roughly 60\% of trace tokens
can be deleted while preserving QA-measured semantic content.

\textbf{Head-level redundancy.} Each attention head maintains its own
copy of the KV cache. Information encoded at a given token is
distributed across heads. When eviction is performed independently per
head (as in Random Attention), the probability that a fact encoded in
$m$ heads is retained by at least one head is:
\[
  P(\text{fact retained}) = 1 - (1-p)^m
\]
For $m = 6$ heads and $p = 0.5$ budget: $P = 1 - 0.5^6 = 0.984$.
This near-certainty of retention means no scored selection is needed.

\section{Technical Formulation}

Let $\mathcal{P}$ be the set of prompt token positions and
$\mathcal{T}$ the set of reasoning-trace positions. Given a cache
budget~$B$, standard evictors compute:
\[
  \mathcal{R} = \mathcal{P} \cup \operatorname{top}_k^{\!\text{score}}(\mathcal{T}),
  \quad k = B - |\mathcal{P}|
\]
Random Attention replaces scored selection with uniform sampling:
\[
  \mathcal{R}^{(h)} = \mathcal{P} \cup
  \operatorname{sample}_k^{\!\text{uniform}}(\mathcal{T})
\]
performed independently per attention head $h$. No score function is
defined or computed.

\section{Experimental Findings}

Experiments span four models --- DeepSeek-R1-Distill-Qwen-7B,
DeepSeek-R1-Distill-Qwen-14B, QwQ-32B, and DeepSeek-R1-32B --- and
six tasks: AIME 2024, AIME 2025, AMC, MATH500, LiveCodeBench, and GPQA
Diamond. At 50\% cache budget:

\begin{center}
\begin{tabular}{lcc}
\toprule
Method & AIME 2024 & MATH500 \\
\midrule
Full KV cache           & 72.3\% & 91.4\% \\
H2O (50\%)              & 68.1\% & 89.3\% \\
StreamingLLM (50\%)     & 67.4\% & 88.7\% \\
\textbf{Random Att.\ (50\%)} & \textbf{68.4\%} & \textbf{89.5\%} \\
\bottomrule
\end{tabular}
\end{center}

Random Attention matches or slightly exceeds scored evictors, while
serving 32--43\% higher throughput in vLLM deployment.

\section{Throughput Analysis}

The throughput advantage has two sources. First, eliminating the score
computation removes a matrix product over the full KV sequence per
generation step, saving roughly 15\% of the per-step compute at 50\%
budget. Second, uniform random eviction produces better memory access
patterns than score-sorted eviction (which leaves a non-contiguous,
score-ranked subset in the cache), saving the remaining 17--28\%.

\section{Ablations and Interpretation}

\textbf{Per-head vs.\ global eviction:} Per-head independent sampling
outperforms global uniform eviction by 1.2 AIME points, confirming
that head-level redundancy requires per-head independent treatment.

\textbf{Tight budgets (25\%):} Random Attention is more robust than
scored evictors at 25\% budget, where its MATH500 accuracy matches
scored evictors at 50\%.

\textbf{Task type:} Temporal and procedural tasks (LiveCodeBench)
benefit slightly more from Random Attention than purely mathematical
tasks, consistent with higher restatement frequency in procedural traces.

\textbf{Prompt fragility confirmed:} If the prompt is evicted even
partially (Condition C), accuracy drops 17+ points on AIME 2024 ---
confirming that the prompt, not the trace, is semantically irreplaceable.

\section*{Reference}

Heng Wang et al. \textbf{Random Attention: Rethinking KV Cache Eviction
for Efficient Reasoning.} arXiv:2609.03430, EMNLP 2026.
\url{https://arxiv.org/abs/2609.03430}

\end{document}
